\documentclass[11pt]{article}
\usepackage[margin=1in]{geometry}
\usepackage{enumitem}
% =============================================================================
% Preambles/header.tex — shared LaTeX/Beamer preamble for this project
%
% Use from a Beamer lecture by prepending:
%     \input{header}
% and compiling with TEXINPUTS=../Preambles:$TEXINPUTS (the /compile-latex
% skill does this automatically).
%
% PALETTE CONTRACT. Color names defined here MUST match the names in
% Quarto/theme-template.scss so Beamer and Quarto renderings use the same
% palette. The scripts/check-palette-sync.sh script enforces this.
%
% To customize: edit the HEX values below AND the matching $variable in
% Quarto/theme-template.scss, then run scripts/check-palette-sync.sh.
% =============================================================================

% -----------------------------------------------------------------------------
% Engine detection + encoding
%
% Our /compile-latex skill uses XeLaTeX, where inputenc is unnecessary and
% can produce encoding warnings. Only load inputenc under pdfTeX.
% -----------------------------------------------------------------------------
\usepackage{iftex}
\ifPDFTeX
  \usepackage[utf8]{inputenc}
\fi

\usepackage{amsmath, amssymb, amsthm}
\usepackage{booktabs}
\usepackage{graphicx}

% -----------------------------------------------------------------------------
% PALETTE — mirrors Quarto/theme-template.scss variable names.
% Load xcolor and define colors BEFORE anything that references them
% (notably hyperref, hypersetup, and the Beamer color assignments).
% -----------------------------------------------------------------------------
\usepackage{xcolor}

\definecolor{primary-blue}{HTML}{012169}      % headings, accents
\definecolor{primary-gold}{HTML}{B9975B}      % emphasis, borders
\definecolor{highlight-yellow}{HTML}{F2A900}  % markers, alerts
\definecolor{light-bg}{HTML}{E8EDF5}          % title / section backgrounds
\definecolor{jet}{HTML}{1A1A1A}               % body text

% Semantic colors (used in diagrams and annotations)
\definecolor{positive}{HTML}{15803D}          % good / observed / identified
\definecolor{negative}{HTML}{B91C1C}          % bad / problematic / bias
\definecolor{neutral}{HTML}{525252}           % reference / context

% Highlight accents (match .hi-* classes in the SCSS)
\definecolor{hi-slate}{HTML}{314F4F}
\definecolor{hi-green}{HTML}{2E8B57}
\definecolor{hi-red}{HTML}{C41E3A}

% Semantic aliases so skill templates and snippets can write
% \textcolor{accent}{...} without hard-coding "primary-blue".
\colorlet{accent}{primary-blue}
\colorlet{accent2}{primary-gold}

% -----------------------------------------------------------------------------
% Hyperref — loaded AFTER xcolor + palette so link colors resolve.
% -----------------------------------------------------------------------------
\usepackage{hyperref}
\hypersetup{colorlinks=true, linkcolor=primary-blue, urlcolor=primary-blue,
            citecolor=primary-blue, pdfborder={0 0 0}}

% -----------------------------------------------------------------------------
% Course code — set once per deck (after \input{header}, before \begin{document})
% via \coursecode{CS401}. Shown in the footer on every slide so a deck's course
% is identifiable without opening the title page. Empty by default.
% -----------------------------------------------------------------------------
\makeatletter
\newcommand{\coursecode}[1]{\gdef\@coursecodetext{#1}}
\gdef\@coursecodetext{}
\makeatother

% -----------------------------------------------------------------------------
% Beamer theme assignments (only apply under Beamer).
%
% We detect Beamer via \@ifclassloaded{beamer}, which is the documented,
% non-internal way. This must be wrapped in \makeatletter/\makeatother
% because \@ifclassloaded contains an @-catcode character.
% -----------------------------------------------------------------------------
\makeatletter
\@ifclassloaded{beamer}{%
  \usefonttheme{professionalfonts}
  \setbeamercolor{structure}{fg=primary-blue}
  \setbeamercolor{frametitle}{fg=primary-blue}
  \setbeamercolor{title}{fg=primary-blue}
  \setbeamercolor{subtitle}{fg=primary-gold}
  \setbeamercolor{author}{fg=primary-blue}
  \setbeamercolor{institute}{fg=neutral}
  \setbeamercolor{date}{fg=primary-gold}
  \setbeamercolor{itemize item}{fg=primary-blue}
  \setbeamercolor{itemize subitem}{fg=primary-gold}
  \setbeamercolor{alerted text}{fg=primary-gold}
  \setbeamercolor{block title}{fg=white, bg=primary-blue}
  \setbeamercolor{block body}{bg=light-bg}
  % Semantic block variants: block=definition/concept (blue), exampleblock=
  % worked example (green/positive), alertblock=key takeaway or caution (gold).
  % Keep to <=2 colored boxes per slide across ALL three variants (INV-7).
  \setbeamercolor{block title example}{fg=white, bg=positive}
  \setbeamercolor{block body example}{bg=light-bg}
  \setbeamercolor{block title alerted}{fg=white, bg=primary-gold}
  \setbeamercolor{block body alerted}{bg=light-bg}

  % Minimal footer: course code (left) + page X / N (right), no navigation clutter
  \setbeamertemplate{navigation symbols}{}
  \setbeamertemplate{footline}{%
    \color{neutral}\footnotesize\hspace{0.7em}\@coursecodetext\hfill
    \insertframenumber/\inserttotalframenumber\hspace{0.7em}\vspace{0.3em}}
}{}
\makeatother

% -----------------------------------------------------------------------------
% TikZ setup — libraries the snippet gallery uses
% -----------------------------------------------------------------------------
\usepackage{tikz}
\usetikzlibrary{arrows.meta, positioning, calc, decorations.pathreplacing,
                fit, shapes.geometric, backgrounds}

% Shared TikZ styles — snippets and hand-written diagrams can pick these up
% via `\begin{tikzpicture}[every node/.style={...}]` or by naming specific
% styles (e.g., `\node[dag-node] (X) at (0,0) {X};`).
\tikzset{
  % Node styles for causal DAGs and similar
  dag-node/.style = {
    draw=primary-blue, fill=light-bg, rounded corners=2pt,
    minimum width=1.2cm, minimum height=0.9cm, align=center, font=\small
  },
  decision-node/.style = {
    draw=primary-gold, fill=white, diamond, aspect=1.8,
    minimum width=1.8cm, minimum height=1.2cm, align=center, font=\small,
    inner sep=1pt
  },
  % Edge styles — observed vs counterfactual
  observed-edge/.style = {-{Latex[length=2.5mm]}, thick, draw=jet},
  counterfactual-edge/.style = {-{Latex[length=2.5mm]}, thick, draw=neutral, dashed},
  confound-edge/.style = {-{Latex[length=2.5mm]}, thick, draw=negative, dashed},
  % Data-point styles for plots
  observed-dot/.style = {fill=primary-blue, draw=primary-blue, circle, inner sep=1.2pt},
  counterfactual-dot/.style = {fill=white, draw=neutral, circle, inner sep=1.2pt}
}

% -----------------------------------------------------------------------------
% Convenience macros
% -----------------------------------------------------------------------------
\newcommand{\muted}[1]{\textcolor{neutral}{#1}}
\newcommand{\key}[1]{\textcolor{primary-gold}{\textbf{#1}}}
\newcommand{\good}[1]{\textcolor{positive}{#1}}
\newcommand{\bad}[1]{\textcolor{negative}{#1}}

% Standout frame macro: full-bleed dark background for section transitions.
% Beamer-only (uses \begin{frame}/\setbeamercolor/\usebeamerfont) -- do not
% call from an article-class file (e.g. Notes/<CODE>/*-notes.tex). \muted,
% \key, \good, \bad, and \coursecode above are plain xcolor/gdef and are
% safe to use from either class; verified when Notes/article-class support
% was added.
% Expand: \transitionslide{Your transition title}
\newcommand{\transitionslide}[1]{%
  \begingroup
  \setbeamercolor{background canvas}{bg=primary-blue}
  \begin{frame}[plain]
    \centering
    \vfill
    {\usebeamerfont{title}\color{white}\Large #1\par}
    \vfill
  \end{frame}
  \endgroup
}

% Section divider frame (Beamer-only), an alternative to \transitionslide with a
% darker two-line style: near-black (jet) background, a small white label line
% above a larger gold title, and a thin gold rule along the bottom edge.
% Expand: \sectiondivider{Section 2}{Pointers}
\newcommand{\sectiondivider}[2]{%
  \begingroup
  \setbeamercolor{background canvas}{bg=jet}
  \begin{frame}[plain]
    \vspace*{3.0cm}
    {\color{white}\ttfamily\large #1\par}
    \vspace{0.5cm}
    {\color{primary-gold}\LARGE #2\par}
    \begin{tikzpicture}[remember picture, overlay]
      \draw[primary-gold, line width=2.5pt]
        ($(current page.south west)+(0,0.1)$) -- ($(current page.south east)+(0,0.1)$);
    \end{tikzpicture}
  \end{frame}
  \endgroup
}

% -----------------------------------------------------------------------------
% Channel watermark — "Concepts That Click" (YouTube).
%
% Article-class files (CompetitiveExam Q/A sets, Notes, Assignments) get a
% light diagonal draft watermark on every page plus a centered footer naming
% the channel. Beamer decks get a subtle TikZ background watermark instead
% (the footline already carries the course code). Centralized here so every
% MCQ PDF is branded without per-file edits.
% -----------------------------------------------------------------------------
\makeatletter
\@ifclassloaded{article}{%
  \usepackage{draftwatermark}
  \SetWatermarkText{Concepts That Click}
  \SetWatermarkScale{1}
  \SetWatermarkFontSize{2cm}
  \SetWatermarkLightness{0.86}
  \SetWatermarkAngle{45}
  \usepackage{fancyhdr}
  \pagestyle{fancy}
  \fancyhf{}
  \fancyfoot[C]{\footnotesize\textcolor{neutral}{Concepts That Click~|~YouTube: \href{https://www.youtube.com/@concepts.thatclick}{@concepts.thatclick}}}
  \renewcommand{\headrulewidth}{0pt}
  \renewcommand{\footrulewidth}{0pt}
  \fancypagestyle{plain}{%
    \fancyhf{}%
    \fancyfoot[C]{\footnotesize\textcolor{neutral}{Concepts That Click~|~YouTube: \href{https://www.youtube.com/@concepts.thatclick}{@concepts.thatclick}}}%
    \renewcommand{\headrulewidth}{0pt}%
    \renewcommand{\footrulewidth}{0pt}%
  }
}{}
\@ifclassloaded{beamer}{%
  \setbeamertemplate{background}{%
    \begin{tikzpicture}[remember picture, overlay]
      \node[rotate=45, opacity=0.055, align=center]
        at (current page.center)
        {\Huge\textcolor{neutral}{Concepts That Click}};
    \end{tikzpicture}%
  }
}{}
\makeatother 
\coursecode{JKSSB}
\renewcommand{\thesection}{27.\arabic{section}}
\newtheorem{example}{Example}[section]
\newenvironment{definitionbox}[1]{%
  \par\noindent\textbf{Definition (#1).}\ \itshape
}{\par\vspace{0.3em}}
\title{Excel Data Tools --- Detailed Lecture Notes}
\author{JKSSB Computer Awareness}
\date{\today}

\begin{document}
\maketitle

\section{What Excel data tools do}
Once data has been entered into a worksheet, it usually has to be organised,
narrowed, summarised, and presented. Excel groups the features that do this
under what we can call its data tools: \textbf{sort}, \textbf{filter},
\textbf{tables}, \textbf{charts}, \textbf{PivotTables and PivotCharts},
\textbf{printing}, \textbf{protection}, and \textbf{what-if analysis}. These are
the features a JKSSB candidate is expected to recognise by name and by
one-line purpose.

\begin{definitionbox}{Data tools}
The worksheet features that arrange, narrow, summarise, present, print, and
safeguard data without changing the underlying values.
\end{definitionbox}

The single idea that holds the whole lesson together is this: most of these
tools change the \emph{view} of the data. Sorting, filtering, and a PivotTable
reorganise what you see on the screen; the original values in the cells stay
exactly where they were. The exceptions are deliberate edits --- a sort does
physically move rows, and a user can type new data --- but the summary tools
themselves are read-only with respect to the source.

\begin{center}
\begin{tabular}{p{0.22\textwidth}p{0.66\textwidth}}
\toprule
\textbf{Tool} & \textbf{One-line purpose} \\
\midrule
Sort & Arrange rows in a chosen order. \\
Filter & Show only the rows that meet a condition. \\
Excel table & A named, auto-expanding, formatted range of data. \\
Chart & A picture of the data, linked to its source range. \\
PivotTable & A fast summary of a large table of data. \\
PivotChart & A chart built on a PivotTable. \\
Protection & Restrict editing of cells or workbook structure. \\
What-if analysis & Test how the result changes with different inputs. \\
\bottomrule
\end{tabular}
\end{center}

\section{Sorting}
\textbf{Sorting} rearranges rows into an order chosen by the user. \textbf{Ascending}
order goes from A to Z, or from the smallest number to the largest.
\textbf{Descending} order goes from Z to A, or from the largest number to the
smallest. A sort is performed on a \emph{column}, but the whole row moves with
that column so that each record stays intact.

The \textbf{Sort} dialog (\key{Data > Sort}) can order data by \textbf{more than
one level}. For example, a payroll sheet can be sorted by Department first, and
then by Salary within each department. The dialog's \textbf{Add Level} button
adds these secondary and tertiary keys.

A \textbf{custom list} sorts by a defined order rather than alphabetically: days
of the week and months are the standard examples. Excel can also sort by cell
colour, font colour, or cell icon. When the range has titles, the option
\textbf{My data has headers} must be ticked so that the header row is not sorted
along with the data.

\begin{example}
A teacher's marksheet has the columns Roll No., Name, and Marks. Clicking
anywhere in the Marks column and choosing Sort Largest to Smallest reorders the
rows so the highest marks appear first. If only the Marks column is selected and
the rest of the row does not move, the records are scrambled --- the names no
longer match the marks.
\end{example}

\textbf{Trap watch.} Sorting a \textbf{single column} on its own scrambles the
records. Always sort the whole table, or convert the data to an Excel table
first.

\section{Filtering}
\textbf{Filtering} displays only the rows that satisfy a condition and hides the
rest. The hidden rows are \textbf{not deleted}; they are temporarily out of
sight. AutoFilter is switched on and off with \key{Data > Filter}, whose
keyboard shortcut is \key{Ctrl}$+$\key{Shift}$+$\key{L}. Once it is on,
drop-down arrows appear in the header row of each column.

\begin{definitionbox}{Filter}
A tool that shows only the rows matching a condition while hiding --- not
removing --- the other rows.
\end{definitionbox}

Clicking a drop-down arrow allows two broad kinds of narrowing:
\begin{itemize}
  \item a \textbf{value filter}, where one or more ticks are kept from a list of
    the column's values;
  \item a \textbf{condition filter}, where a rule is applied. Text Filters,
    Number Filters, and Date Filters offer common rules such as \textit{Equals},
    \textit{Greater Than}, \textit{Begins With}, and \textit{Between}.
\end{itemize}
\textbf{Top 10} and \textbf{Above Average} are ready-made condition filters.
\textbf{Advanced Filter} is used for more complex criteria and can copy the
matching rows to another location on the sheet.

Clearing the filter with \textbf{Clear Filter} restores every row. The key
contrast to remember is that filtering \emph{hides} rows; it does not remove
them and does not change the data.

\begin{example}
A sales sheet has a Region column. Applying a value filter and ticking only
``North'' hides all other regions. The rows for South, East, and West are still
present in the worksheet; clearing the filter brings them back unchanged.
\end{example}

\section{Excel tables}
An \textbf{Excel table} is a range of data that has been converted with
\key{Insert > Table} (shortcut \key{Ctrl}$+$\key{T}) and given a table style. A
table has a \textbf{name} (by default \textit{Table1}), and it behaves as a
single object rather than as loose cells.

\begin{definitionbox}{Excel table}
A named, formatted range of data that automatically expands to include new
rows and columns and whose formulas use structured references.
\end{definitionbox}

The practical advantages of a table over a plain range are:
\begin{itemize}
  \item it \textbf{auto-expands} when a new row is typed immediately below it;
  \item it shows \textbf{banded rows} that improve readability;
  \item it adds \textbf{filter buttons} to the header row automatically;
  \item it can display an optional \textbf{Total Row}; and
  \item its formulas use \textbf{structured references} such as
    \key{Table1[Sales]} instead of cell addresses such as \key{B2}.
\end{itemize}
Because a structured reference points at the whole column, it automatically
covers new rows that are added to the table. A formula written with cell
addresses would have to be extended by hand.

\begin{center}
\begin{tabular}{lll}
\toprule
\textbf{Point} & \textbf{Plain range} & \textbf{Excel table} \\
\midrule
Name & None & Assigned name (Table1) \\
Grows & Manually & Auto-expands \\
Formulas & Cell addresses & Structured references \\
Filters & Manual AutoFilter & Built-in header buttons \\
\bottomrule
\end{tabular}
\end{center}

A table is \emph{not} a PivotTable. A table is the \textbf{source data}; a
PivotTable is one of the summaries that can be built on top of it.

\section{Charts}
A \textbf{chart} is a graphical picture of worksheet data. Because the chart is
\textbf{linked} to its source range, editing the numbers in the cells updates
the chart automatically. The common chart types and their uses are:

\begin{center}
\begin{tabular}{p{0.26\textwidth}p{0.62\textwidth}}
\toprule
\textbf{Chart type} & \textbf{Best used for} \\
\midrule
Column / bar & Comparing values across categories. \\
Line & Showing a trend, usually over time. \\
Pie / doughnut & Showing the parts of a single whole. \\
Scatter (XY) & Showing the relationship between two numeric variables. \\
\bottomrule
\end{tabular}
\end{center}

The parts of a chart carry their own names, and exam items often test them. The
\textbf{chart area} is the whole chart, including the frame. Inside it, the
\textbf{plot area} is the region where the data is actually drawn. Each set of
values is a \textbf{data series}, and the \textbf{legend} names each series. The
\textbf{category axis} (usually the X-axis) carries the labels, while the
\textbf{value axis} (the Y-axis) carries the numbers. \textbf{Data labels}
display the exact value at each data point, and \textbf{gridlines} help read the
values.

A chart normally floats on the same worksheet as its data, but it can be moved
to a separate \textbf{chart sheet} that contains only the chart.

\begin{example}
A school records the number of students in each class. A column chart compares
the classes at a glance. The same data as a line chart would suggest a trend
across the classes, which may not be meaningful; for parts of one total, such as
the share of boys and girls in the school, a pie chart is the right choice.
\end{example}

\section{PivotTables and PivotCharts}
A \textbf{PivotTable} summarises a large table of data in seconds. Instead of
writing formulas, the user \textbf{drags} fields into four areas:

\begin{center}
\begin{tabular}{p{0.20\textwidth}p{0.66\textwidth}}
\toprule
\textbf{Area} & \textbf{What goes there} \\
\midrule
Rows & Field values listed down the left side of the summary. \\
Columns & Field values spread across the top of the summary. \\
Values & The numbers that are summarised (Sum, Count, Average, and more). \\
Filters & Fields used to filter the whole PivotTable. \\
\bottomrule
\end{tabular}
\end{center}

The \textbf{Values} area chooses the aggregation: by default a numeric field is
summed and a text field is counted, but Count, Average, Max, Min, and other
functions can be chosen. Rearranging the field layout re-computes the summary
instantly --- this quick re-orientation is what gives the PivotTable its name.

\begin{definitionbox}{PivotTable}
An interactive summary that aggregates a large range of data by dragging fields
into Rows, Columns, Values, and Filters areas.
\end{definitionbox}

A \textbf{PivotChart} is a chart built on a PivotTable and linked to it. When
the PivotTable layout changes, the PivotChart changes with it.

The most important contrast is that a PivotTable \textbf{summarises} the source
data; it does not replace it and does not change it. If the source data is
edited, the PivotTable must be \textbf{refreshed} to show the new result.
Deleting the source range breaks the PivotTable, and a PivotTable is never the
same thing as an Excel table.

\section{Printing worksheets}
A worksheet is printed from \key{File > Print}, and its layout is controlled
from the \key{Page Layout} tab. The key settings are:
\begin{itemize}
  \item \textbf{Orientation} --- \textbf{Portrait} for a tall page or
    \textbf{Landscape} for a wide one;
  \item \textbf{Paper size}, \textbf{margins}, and \textbf{headers/footers};
  \item \textbf{Print Area} --- select a range and choose \key{Set Print Area}
    to print only that part of the sheet;
  \item \textbf{Print Titles} --- repeats chosen header rows or columns at the
    top or left of every printed page, so a long table keeps its headings; and
  \item \textbf{Scaling} --- options such as ``Fit Sheet on One Page'' shrink
    the printout to a chosen width or number of pages.
\end{itemize}
\textbf{Gridlines} are the faint cell borders seen on screen. They are
\textbf{not printed by default}; the \key{Print} gridlines option under Sheet
Options must be enabled to show them on paper. \textbf{Page breaks} can be
inserted to control where a new printed page begins.

\begin{example}
A 200-row price list is being printed. Without Print Titles, only the first page
shows the column headings and later pages are unreadable. Setting the heading
row as a print title repeats it on every page. If the list is a little too wide,
landscape orientation or ``Fit All Columns on One Page'' keeps it tidy.
\end{example}

\section{Protecting worksheets and workbooks}
\textbf{Protection} restricts changes that can be made to a workbook.
\textbf{Protect Sheet} (\key{Review > Protect Sheet}) prevents cells on that
sheet from being edited, while allowing selected actions such as selecting
locked cells or formatting. \textbf{Protect Workbook} protects the
\textbf{structure} of the workbook --- preventing sheets from being added,
deleted, renamed, or moved --- and can also protect its windows.

\begin{definitionbox}{Protection}
A restriction placed on a worksheet or workbook that prevents specified changes,
optionally guarded by a password.
\end{definitionbox}

Cell locking is the detail most candidates miss. Every cell is \textbf{Locked}
by default in the Format Cells dialog, but this lock has \textbf{no effect until
the sheet is protected}. To let a user edit some cells and not others, the
chosen cells must first be \textbf{unlocked}, and then the sheet must be
protected.

Protection is not the same as encryption. Protection restricts editing inside
Excel; it is not the password that encrypts the whole file (that is a separate
file-level password).

\section{What-if analysis}
\textbf{What-if analysis} answers the question ``what would the result be if an
input changed?''. Excel provides three tools for this at the depth a JKSSB
candidate needs:

\begin{center}
\begin{tabular}{p{0.26\textwidth}p{0.62\textwidth}}
\toprule
\textbf{Tool} & \textbf{What it does} \\
\midrule
Goal Seek & Finds the input value needed to reach a target result. \\
Scenario Manager & Stores and compares several sets of input values. \\
Data Table & Shows how an output changes as one or two inputs change. \\
\bottomrule
\end{tabular}
\end{center}

All three are found under \key{Data > What-If Analysis}. Goal Seek works
backwards from a target: for example, it finds the selling price that makes the
profit exactly a desired figure. Scenario Manager keeps named alternatives such
as ``Best case'' and ``Worst case''. A Data Table is a grid that shows a range
of results for one or two changing inputs.

\section{Near-neighbour contrast: COUNTIFS vs SUMPRODUCT}
To \textbf{count rows that satisfy several conditions at the same time}, Excel
offers two tools that students often confuse.

\textbf{COUNTIFS} is the direct multi-criterion counting function. Its arguments
are written as \textbf{paired} ranges and criteria:
\[
\texttt{COUNTIFS(range1, crit1, range2, crit2, \ldots)}
\]
It applies \textbf{AND} logic: a row is counted only if it meets \emph{every}
range-and-criteria pair. For example,
\texttt{COUNTIFS(A2:A100,"North",B2:B100,">500")} counts the rows that are in
the North region \emph{and} have a value greater than 500.

\textbf{SUMPRODUCT} is a general array function. It can also count
multi-criterion rows by multiplying boolean results:
\[
\texttt{SUMPRODUCT((A2:A100="North")*(B2:B100>500))}
\]
Each TRUE/FALSE comparison becomes 1 or 0 when multiplied, and SUMPRODUCT adds
the products, giving the same count. Because the multiplication already coerces
the boolean values to numbers, SUMPRODUCT does not need the \key{--} (double
unary) that is sometimes seen when boolean arrays are fed to other functions.

For \textbf{OR} logic --- a row that meets either of two conditions --- the two
matches are \emph{added}, for example
\texttt{COUNTIF(range,"North")+COUNTIF(range,"South")}.

\begin{example}
A marksheet has a Stream column and a Marks column. To count students who are in
the Science stream \emph{and} scored more than 80, write
\[ \texttt{COUNTIFS(Stream,"Science",Marks,">80")}. \]
The same count can be obtained with
\[ \texttt{SUMPRODUCT((Stream="Science")*(Marks>80))}. \]
Confusing the paired arguments of COUNTIFS, or thinking COUNTIFS cannot handle
two conditions at once, is the classic error (TRAP-14, PYQ-15).
\end{example}

\section{Exam retrieval and common confusions}
Six distinctions prevent most errors on this topic.
\begin{enumerate}
  \item \textbf{Sort} rearranges rows into an order; \textbf{filter} shows only
    matching rows and \textbf{hides} the rest --- it does not delete them.
  \item A \textbf{PivotTable} summarises the source data but is \textbf{not} the
    source table; it must be \textbf{refreshed} after the data changes.
  \item A PivotTable is \textbf{not} an Excel table. The table is the source;
    the PivotTable is a summary built on it.
  \item \textbf{COUNTIFS} counts with AND across \textbf{paired}
    range-and-criteria arguments; \textbf{SUMPRODUCT} does the same with
    boolean arrays (TRAP-14).
  \item \textbf{Locked} cells are protected only \textbf{after} the sheet is
    protected; the default lock has no effect on its own.
  \item \textbf{Gridlines} are \textbf{not printed} unless Print gridlines is
    enabled.
\end{enumerate}

\begin{example}
An item states, ``A PivotTable changes the source data and updates on its own.''
Both halves are wrong. A PivotTable never changes the source; and after the
source changes it does not update by itself --- the user must refresh it. The
stable concept is: PivotTable = summary, source = untouched, refresh = manual.
\end{example}

\subsection*{Exercises}
\begin{enumerate}[label=\arabic*.]
  \item Distinguish between sorting and filtering, and state what happens to the
    rows that a filter does not show.
  \item Name two features that an Excel table provides but a plain range does
    not.
  \item Match each chart to its best use: column, line, pie, scatter.
  \item List the four areas of a PivotTable and say what goes into each.
  \item Explain why a PivotTable must be refreshed after the source data
    changes.
  \item What is the difference between Protect Sheet and Protect Workbook?
    When does a locked cell actually become protected?
  \item Write the COUNTIFS formula that counts rows in the North region with a
    value greater than 500, and give the equivalent SUMPRODUCT formula.
\end{enumerate}

\subsection*{Solutions}
\begingroup\small
\begin{enumerate}[label=\arabic*.]
  \item Sorting rearranges rows into a chosen order; filtering shows only the
    rows matching a condition. A filter \emph{hides} the other rows temporarily;
    it does not delete them, and clearing the filter restores them.
  \item An Excel table has a name and auto-expands when new rows are typed; it
    also provides structured references, banded rows, and built-in filter
    buttons. A plain range has none of these by default.
  \item A column chart compares values across categories; a line chart shows a
    trend over time; a pie chart shows parts of one whole; a scatter chart shows
    the relationship between two numeric variables.
  \item Rows --- field values down the left; Columns --- field values across the
    top; Values --- the numbers that are summarised; Filters --- fields used to
    filter the whole PivotTable.
  \item A PivotTable stores only a summary view; it does not track the source
    automatically. After the source data changes, the user must Refresh the
    PivotTable so the summary is recalculated.
  \item Protect Sheet stops cells on a sheet from being edited; Protect Workbook
    protects the workbook structure (adding, deleting, renaming, or moving
    sheets). A cell is Locked by default, but the lock takes effect only after
    the sheet is protected.
  \item \texttt{COUNTIFS(A2:A100,"North",} \texttt{B2:B100,">500")}; the
    equivalent is \texttt{SUMPRODUCT((A2:A100="North")}
    \texttt{*(B2:B100>500))}.
\end{enumerate}
\endgroup
\end{document}
